\unit{第212次报告——商预概形}
\sid{212}
\src{99}
\seminarzh{第13年度，1960/61，第212号\\（1961年2月）}
\paperhead{构造技术与代数几何中的存在性定理\\
III. 商预概形}%
  {亚历山大·格罗滕迪克 著}

\subsection*{引言}
\addcontentsline{toc}{subsection}{引言}

本讲所处理的问题与前两讲所考察的问题不同，因为这里试图表示的是以变动概形为
变量的某些协变函子，而不再是反变函子。不过，取商过程在代数几何的许多构造问题中
都是本质性的，其中也包括讲稿 I 和 II 中的问题（[1]、[2]）。例如，在 $T$-预概形
$X$ 上，相对于忠实平坦且拟紧的态射 $T\longrightarrow S$，一项\emph{下降数据的有效性}
问题，等价于如下问题：由该下降数据在 $X$ 中定义的平坦等价关系是否存在一个
$X$ 的商（并满足下文考察的合理性质）；[1]，A，\S~2 c 中提出的问题，大概会与
本讲第 2 节提出的问题同时得到解决。同样，$S$-概形 $X$ 的
\emph{Picard 概形}（定义见 [2]，C，\S~3）可以用多种方式定义为某些其他概形
（正因子概形，或到某个射影概形中的浸入概形）关于平坦等价关系的商；而这些辅助
概形的定义和构造在技术上更为简单：事实上，它们是 [2]，C，\S~2 中定义的
$\underline{\operatorname{Hom}}_S(X,Y)$ 型概形及其变体；在适当的射影性假设下，
它们的构造将是下一讲的主题。因此，结合本讲与下一讲的结果，便可在适当假设下构造
Picard 概形。

在预概形范畴中取商的问题仍包含若干尚未解决的疑问。其中最重要的一个见第 8 节。
目前，它仍是构造\emph{整数环上任意亏格曲线、极化阿贝尔簇等的模概形}的唯一障碍。
由此可见，这个问题值得代数群专家为之努力。

\src{100}
\subsection*{1. 等价关系与有效等价关系}
\addcontentsline{toc}{subsection}{1. 等价关系与有效等价关系}

设 $\mathcal C$ 是一个范畴，$X$ 是 $\mathcal C$ 的一个对象。若一对态射
\[
  p_1,p_2:R\rightrightarrows X
\]
满足如下条件，就称它为 $\mathcal C$ 中一个\emph{等价对}，其\emph{靶}为 $X$，
\emph{源}为 $R$：对 $\mathcal C$ 的每个对象 $T$，相应的两个映射
\[
  p_1(T),p_2(T):R(T)\rightrightarrows X(T)
\]
（这里对 $\mathcal C$ 的每个对象 $Y$，约定
$Y(T)=\operatorname{Hom}(T,Y)$）定义一个映射
\[
  R(T)\longrightarrow X(T)\times X(T)
\]
它把 $R(T)$ 双射地映到集合 $X(T)$ 上某个等价关系的图。靶为 $X$ 的各个等价对
之间有一个显然的等价关系；关于后一等价关系的一个等价类称为 $X$ 中的一个
\emph{$\mathcal C$-等价关系}；不会引起混淆时，简称等价关系。

若 $X\times X$ 存在，则在 $X$ 中给定一个等价关系，等价于给定
$X\times X$ 的一个子对象 $R$，使得对 $\mathcal C$ 的每个对象 $T$，在
\[
  (X\times X)(T)=X(T)\times X(T)
\]
中与 $R(T)$ 相应的子集，是 $X(T)$ 上某个等价关系的图。以 $p_1$ 和 $p_2$
表示由投影 $\operatorname{pr}_1$ 和 $\operatorname{pr}_2$ 诱导的从 $R$
到 $X$ 的态射，则上述条件正表示 $(p_1,p_2)$ 是一个等价对。在
$X\times X$ 以及纤维积 $(R,p_2)\times_X(R,p_1)$ 存在的前提下，也可以按照
[2]，A，\S~1 中说明的一般原则，在 $\mathcal C$ 中用交换图表示 $X(T)$ 中
$R(T)$ 所满足的集合论等价关系公理。我们不需要采用这种表述。

每当给定一对具有共同源 $R$ 和共同靶 $X$ 的态射 $(p_1,p_2)$ 时，可以把这一对的
\emph{余核}定义为 $\mathcal C$ 中表示如下关于 $Z$ 协变的函子
\[
  \operatorname{Hom}_{p_1,p_2}(X,Z),
\]
的对象 $Y$；该函子的值是所有态射 $u$（从 $X$ 到 $Z$）中满足
$up_1=up_2$ 者所成的集合。若 $Y$ 存在，则它在唯一同构意义下确定。将它记为
$X/(p_1,p_2)$，或滥用记号记为 $X/R$；后一记号主要用于
\src{101}$(p_1,p_2)$ 是等价对的情形：此时习惯上在记号中把该对定义的等价关系
与 $R$ 视为同一对象。注意，若把 $Y$ 看作 $X$ 的商，则它事实上只依赖于
$(p_1,p_2)$ 所定义的等价关系。

现在从一个态射
\[
  f:X\longrightarrow Y
\]
出发；它使我们可以把 $X$ 看作一个“$Y$ 上的对象”。假设纤维积
\[
  R(f)=X\times_Y X
\]
存在，并以 $p_1$、$p_2$ 表示它的两个投影。于是 $(p_1,p_2)$ 是一个等价对，
称为与态射 $f$ \emph{相伴的}等价对。因此，它定义一个等价关系，称为与 $f$
相伴的等价关系。

设 $(p_1,p_2)$ 是一对靶为 $X$、共同源为 $R$ 的态射。若
\begin{enumerate}
  \item[(i)] 余核 $X/(p_1,p_2)=Y$ 存在；
  \item[(ii)] 纤维积 $X\times_Y X$ 存在；
  \item[(iii)] 以 $p_1$、$p_2$ 为分量的态射
  $R\longrightarrow X\times_Y X$ 是同构，
\end{enumerate}
就称 $(p_1,p_2)$ 为一个\emph{有效等价对}。此时 $(p_1,p_2)$ 确实是一个
等价对。它所定义的等价关系也称为\emph{有效等价关系}。

若态射 $f:X\longrightarrow Y$ 满足
\begin{enumerate}
  \item[(i)] 纤维积 $X\times_Y X=R$ 存在；
  \item[(ii)] 商 $X/(p_1,p_2)$ 存在，其中 $p_1$、$p_2$ 是从 $R$ 到 $X$
  的两个投影；
  \item[(iii)] 由 $f$ 诱导的态射
  $X/(p_1,p_2)\longrightarrow Y$ 是同构，
\end{enumerate}
就称 $f$ 为\emph{有效满态射}。此时 $f$ 确实是满态射，甚至是严格满态射
（见 [1]，A，\S~2.3）；若纤维积 $X\times_Y X$ 存在，则逆命题也成立。
由满态射 $f$ 定义的 $X$ 的商对象也称为 $X$ 的一个\emph{有效商}。

上述定义蕴含如下“Galois 对应”：

\src{102}
\result{命题}{1.1} $X$ 中的有效等价关系 $R$ 所成的集合，与 $X$ 的有效商 $Y$
所成的集合之间，存在一个保持自然偏序的双射对应：$R$ 对应于有效商 $X/R$；
$Y$ 对应于由典范投影 $X\longrightarrow Y$ 定义的有效等价关系
（该关系由带有两个投影的纤维积 $X\times_Y X$ 定义）。

在性质极好的范畴中（集合、集合层等），每个商都是有效的，每个等价关系也都是
有效的。但在给定预概形 $S$ 上的预概形范畴等范畴中，情况不再如此；即使 $S$
是一个域的谱，甚至即使只考虑在 $S$ 上有限的概形，也是这样。有效性问题，
乃至商的存在性问题（对于在 $S$ 上不有限的预概形），通常都相当棘手。

\subsection*{2. 例：$S$ 上有限的预概形}
\addcontentsline{toc}{subsection}{2. 例：S 上有限的预概形}

设 $\mathcal C$ 是在 $S$ 上有限的预概形所成的范畴，并假设 $S$ 局部
诺特。它等价于 $S$ 上凝聚交换代数层所成范畴的反范畴；或者，当 $S$ 是环
$A$ 的仿射概形时，等价于在 $A$ 上有限的 $A$-代数所成范畴的反范畴
（即这些代数作为 $A$-模是有限生成的）。立即可知，$\mathcal C$ 中有限逆极限
与有限直极限都存在。对前者而言，这是熟知的事实（而且不需要任何有限性假设
\ldots）。例如，$S$ 上预概形 $X,Y$ 的纤维积，对应于相应代数的张量积
$B\otimes_A C$；由两个 $A$-代数同态
$u,v:C\rightrightarrows B$ 所定义的两个态射 $X\rightrightarrows Y$ 的核，
对应于 $B$ 关于由所有 $u(c)-v(c)$ 生成的理想的商，等等。对于有限直极限，
只须一方面考察有限和——它们对应于 $A$-代数的通常直积——另一方面考察态射对
$X\rightrightarrows Y$ 的余核；后者事实上对应于 $C$ 中由所有使同态
$u,v:C\rightrightarrows B$ 取相同值的元素组成的子环（这一点立即可见，而由
诺特假设，该子环在 $A$ 上有限）。还可以利用诺特假设证明：在
$S$ 上有限的预概形范畴 $\mathcal C$ 中如此构造的有限直极限，特别是这些商，
事实上也是\emph{所有}预概形所成范畴中的商。

\src{103}
如同我们在 [1] 中所说，$\mathcal C$ 中存在非有效满态射
（或非严格满态射\footnote{印刷文本在“严格”之前漏掉了“非”。在这里考察的范畴中，
纤维积存在，而上一段恰好证明了有效性与严格满态射性等价。}；因为纤维积存在，
这两种说法等价）。在不作平坦性假设时，\emph{我不知道等价关系是否总是有效的}。
在这个方向上，我只得到了一些很不完全的正面结果；这些结果对于证明形式模理论中的
基本定理是不可缺少的（见 [2]，B，定理 1）。应当指出，对所提出的问题，很容易
约化到如下情形：$S$ 是一个阿廷局部环 $A$ 的谱，而 $A$ 的剩余域代数闭。
但即使 $A$ 是一个域，答案仍然未知。

也可以考察 $S$ 上一个不再假设在 $S$ 上有限的预概形 $X$，同时考虑 $X$ 中一个
等价关系 $R$，使得 $p_1:R\longrightarrow X$ 是有限态射。这时称 $R$ 为
有限等价关系。为简单起见，假设 $S$ 和 $X$ 都是仿射的（于是 $R$ 也是仿射的，
从而情形归结为纯粹交换代数中的情形），则\emph{即使在这种情况下，我们也不知道
商 $X/R=Y$ 是否存在，以及典范态射 $X\longrightarrow Y$ 是否有限}。
（最简单的情形是：$S$ 是一个域 $k$ 的谱，而 $X$ 是 $k[t]$ 的谱，即仿射直线。）
当然，若前述两个问题都有肯定答案，则在当前情形中可断定 $R$ 是有效的。注意，
若把 $X$ 中等价关系的图换成第 4 节意义下 $X$ 中的一个等价预图，则商 $Y$ 的
\emph{存在性}以及 $f:X\longrightarrow Y$ 的有限性问题，仍以完全相同的形式
出现。

通过某个给定的闭子预概形将若干预概形 $X_i$ “粘合”起来以构造预概形时，会出现
关于任意程度的有限等价关系取商的问题；粘合律恰好由各 $X_i$ 的和预概形 $X$ 中
一个有限等价关系表示。还应预期，要沿 [2] 所开创的方向建立一般的非射影构造技术，
首先必须解决这里提出的问题及其各种变体。

撰稿人所知唯一的一般性正面结果如下：

\result{命题}{2.1} 设 $S$ 是局部诺特预概形，$s$ 是 $S$ 的一个点，
$\Omega$ 是 $k(s)$ 的一个代数闭扩域。考虑
\src{104}相应的“纤维函子” $F$：它把每个在 $S$ 上有限的 $S$-概形 $X$
映到 $X_s$ 的所有 $\Omega$-值点所成的集合。这个函子（显然左正合）
\emph{右正合}，即它与有限直极限交换，特别是与态射对的余核交换。

把这一结果用于 $S$ 的所有“几何点”，可推出 $\mathcal C$ 的“商”范畴
$\mathcal C'$——它通过“模去满射纯不可分态射”进行推理而得到
（即形式地为所讨论的态射添加逆）——是一个“几何”范畴，也就是说，它满足与
集合范畴相同的“有限性质”。特别地，其中每个等价关系都是有效的。这意味着，
若 $R$ 是在 $S$ 上有限的 $X$ 中的一个等价关系，则典范态射
\[
  R\longrightarrow X\times_Y X\qquad (\text{其中 }Y=X/R)
\]
是\emph{纯不可分且满射的}（事实上，因为它是单态射，所以它是一个满射闭浸入）。

\subsection*{3. 作用群的情形}
\addcontentsline{toc}{subsection}{3. 作用群的情形}

再次假设 $\mathcal C$ 是任意范畴。设 $G,X$ 是 $\mathcal C$ 的对象，并假设
$G$ 是作用于 $\mathcal C$ 的对象 $X$ 的一个 $\mathcal C$-群。这意味着
（见 [2]，A，\S~1），对 $\mathcal C$ 的每个对象 $T$，在 $G(T)$ 上给定了
群结构，并在 $X(T)$ 上给定了带 $G(T)$ 作用的集合结构，而且当 $T$ 变动时，
这些结构随 $T$ “函子性地变动”。若乘积 $G\times G$ 和 $G\times X$ 在
$\mathcal C$ 中存在，则这种结构也可以定义为一对态射
\[
  G\times G\longrightarrow G,
  \qquad \pi:G\times X\longrightarrow X,
\]
并要求对 $\mathcal C$ 的每个对象 $T$，集合 $G(T)$ 与 $X(T)$ 上相应的合成律
使 $G(T)$ 成为一个作用于 $X(T)$ 的群。把这条公理翻译成 $\mathcal C$ 中某些
图表的交换性并不困难，却很繁琐；事实上，在我所知的一切情形中，这样做都完全
没有必要。

假设 $G\times X$ 存在，并考虑两个态射
\[
  p_1,p_2:G\times X\rightrightarrows X
\]
其中
\src{105}
\[
  p_1=\operatorname{pr}_2,\qquad p_2=\pi.
\]
\footnote{原印本及法、英可编辑见证均作
$p_1=\operatorname{pr}_1$；但
$\operatorname{pr}_1:G\times X\to G$，而这里两个态射的靶都是 $X$，
且下一个公式要求 $(g,x)\mapsto(x,gx)$，故此处采用类型正确的第二投影。}
立即可见，$(p_1,p_2)$ 是等价对，当且仅当对 $\mathcal C$ 的每个对象 $T$，
由这一对定义的映射
\[
  G(T)\times X(T)\simeq(G\times X)(T)
  \longrightarrow X(T)\times X(T),
\]
是单射；也就是说，当且仅当群 $G(T)$ \emph{自由地}作用于集合 $X(T)$；
换言之，若 $g\in G(T)$、$x\in X(T)$ 且 $g\cdot x=x$，则 $g$ 必为群
$G(T)$ 的单位元。这时称 $G$ \emph{自由作用于 $X$}
（或者称 $X$ 是一个\emph{$G$ 下的 $\mathcal C$-主空间}）。这时，与
$(p_1,p_2)$ 相伴的等价关系称为\emph{由自由作用于 $X$ 的群 $G$ 所定义的
等价关系}。若 $X\times X$ 也存在，并考虑由对 $(p_1,p_2)$ 定义的态射
\[
  p:G\times X\longrightarrow X\times X
\]
则 $G$ 自由作用这一条件表示 $p$ 是\emph{单态射}。

当然，即使 $G$ 在 $X$ 上的作用并不自由，我们仍希望得到 $X$ 关于 $G$ 的商
存在的判据，也就是上述对 $(p_1,p_2)$ 的余核存在的判据。

所讨论的余核通常记为 $X/G$；若 $G$ 从左作用，则最好记为 $G\backslash X$
（前一记号留给 $G$ 从右作用的情形）。注意，即使 $G\times X$ 在 $p$ 下的
“像”存在——例如，把这个像定义为 $X\times X$ 中使 $p$ 可以通过它分解的最小
子对象——并将其记为 $R$，它通常也不是 $X$ 中的等价关系。如果这时试图直接
关于 $R$ 取商（更确切地说，关于由两个投影 $\operatorname{pr}_i$ 诱导的从
$R$ 到 $X$ 的态射对取商），就会不可挽回地丢失原始对 $(p_1,p_2)$ 的特殊性质。
因此，必须找到等价关系概念的一种推广，使它能够直接应用于由一个
$\mathcal C$-作用群定义的态射对。

\subsection*{4. 预等价关系}
\addcontentsline{toc}{subsection}{4. 预等价关系}

回顾：若一个范畴的所有态射都是同构，就称它为\emph{群胚}。另一方面，一个范畴
应当定义为由两个基础集合 $(X,R)$ 组成：\emph{对象}集合与\emph{箭头}集合；
并赋予如下结构：
\src{106}
\begin{enumerate}
\item[(i)] 一对映射
\[
  p_1,p_2:R\rightrightarrows X
\]
分别称为\emph{源映射}和\emph{靶映射}，

\item[(ii)] 一个映射
\[
  \pi:(R,p_2)\times_X(R,p_1)\longrightarrow R
\]
称为\emph{复合映射}。
\end{enumerate}
这些数据必须满足一些众所周知的公理，这里不再复述；这些公理可以明确表述为
若干图的交换性，以及存在一个映射 $D:X\to R$（它必然唯一），使另外两个图
交换。$D$ 对应于把一个对象送到相应恒等箭头的运算，并满足
\[
  p_1D=p_2D=\operatorname{id}_X.
\]
说这个范畴是群胚，是指存在一个映射
\[
  s:R\longrightarrow R
\]
（它必然唯一），称为 $R$ 的\emph{对称映射}，它把每支箭头变为一支逆箭头。
这一条件可以明确表述为由 $s$、$\Delta$ 和前述数据构成的另外四个图的
交换性，其中前两个写作
\[
  p_1s=p_2,\qquad p_2s=p_1.
\]

回顾以上各点后，([2]，A，\S~1) 中的一般定义特别说明了，应当如何理解任意
范畴 $\mathcal C$ 的一对对象 $X,R$ 上的
$\mathcal C$-\emph{范畴}结构，或相应的 $\mathcal C$-\emph{群胚}结构：
按照定义，这是对 $\mathcal C$ 的每个对象 $T$ 给出一个集合论意义下的范畴
结构或群胚结构，其对象集为 $X(T)$、箭头集为 $R(T)$，而且这些结构随变化的
$T$ “以函子方式变化”。因此，这些数据给出
两个态射
\[
  p_1,p_2:R\rightrightarrows X
\]
分别称为\emph{源态射}和\emph{靶态射}；当所需的纤维积存在时，还给出一个态射
\src{107}
\[
  \pi:(R,p_2)\times_X(R,p_1)\longrightarrow R
\]
称为\emph{复合态射}。这三项数据足以确定 $X,R$ 上的范畴结构；须对它们施加
的公理如下：对每个 $T$，在 $X(T),R(T)$ 上所得的相应三项数据，都在这对
集合上定义一个范畴结构（或群胚结构）。适当时，这可以表述为若干图的交换性；
这些图涉及一个唯一确定的态射
\[
  D:X\longrightarrow R
\]
以及在群胚情形下一个唯一确定的态射
\[
  s:R\longrightarrow R,
\]
而这些图可完全按照“集合论”情形写出。所幸在实际使用中，无须对这些公理作
如此繁琐的解释。能够借助某些纤维积之间的态射及其等式来表述上述数据与公理，
其唯一的理论意义如下：若函子 $F:\mathcal C\to\mathcal C'$ 左正合（即与
有限积和纤维积交换），它就把一个 $\mathcal C$-范畴（或
$\mathcal C$-群胚）变为一个 $\mathcal C'$-范畴（或
$\mathcal C'$-群胚）（这里要求 $\mathcal C$ 中存在有限积和纤维积）。

在实际使用中，重要的是把态射 $p_1,p_2,\pi,D,s$ 解释为 $\mathcal C$ 中某个
适当半单纯对象或单纯对象的\emph{单纯运算}，至少在 $\mathcal C$ 中存在
纤维积时如此。为固定术语，引入\emph{标准单形范畴 $\mathcal S$}：其对象是
形如
\[
  \Delta_n=[0,n],\qquad\text{其中 }n\geq 0
\]
的有限集合（即从 $0$ 到 $n$ 的整数区间），其态射或箭头则是这些有限集合
之间的\emph{任意映射}。注意，范畴 $\mathcal S$ 等价于以有限集合之间的
映射为态射的\emph{非空}有限集合范畴。在 $\mathcal S$ 中，任意有限
\emph{非空}对象族的余积显然存在，两个对象在第三个对象之下的推出也存在
（这是纤维积的对偶运算）。以 $\mathcal S'$ 表示 $\mathcal S$ 的如下
子范畴：它具有相同的对象，但其态射仅为各 $\Delta_n$ 之间的
\emph{保序映射}。这个范畴等价于非空有限全序集的范畴。在这个范畴中，
\src{108}
两个对象的余积从不存在，而两个对象 $A,B$ 在第三个对象 $C$ 之下的推出
一般也不存在（例如取 $C=\Delta_0$、$A=B=\Delta_1$，并令两个结构映射
$u:C\to A$ 和 $v:C\to B$ 相同）。不过，它在某些情形下确实存在，例如
\[
  C=\Delta_0,\qquad A=\Delta_m,\qquad B=\Delta_n,\qquad
  u(0)=m,\qquad v(0)=0,
\]
此时有
\[
  \Delta_m\amalg_{\Delta_0}\Delta_n=\Delta_{m+n}.
\]

范畴 $\mathcal C$ 中的一个\emph{单纯对象}（相应地，一个
\emph{半单纯对象}），按照定义，是从 $\mathcal S$（相应地，从
$\mathcal S'$）到 $\mathcal C$ 的反变函子 $K$。因此，单纯对象经限制给出
半单纯对象；但前一个概念与后一个概念的本质区别，在于
$K_n=K(\Delta_n)$ 上存在\emph{对称运算}，它们对应于 $n+1$ 个字母上的
对称群各元素在函子 $K$ 下的像（该对称群视为 $\Delta_n$ 在
$\mathcal S$ 中的自同构群）。

现在，对每个 $n$，令 $\Delta'_n$（相应地，$\Delta''_n$）为如下有限范畴：
其对象集为 $\Delta_n$，并由 $\Delta_n$ 上的混沌序关系（相应地，自然全序
关系）定义（箭头集即所述序关系的图）。显然，$\Delta'_n$（相应地，
$\Delta''_n$）以函子方式依赖于 $\mathcal S$（相应地，$\mathcal S'$）的
对象 $\Delta_n$。若 $Z$ 是一个范畴，则当 $\Delta_n$ 变化时，
$\operatorname{Hom}(\Delta'_n,Z)$（相应地，
$\operatorname{Hom}(\Delta''_n,Z)$）给出从范畴 $\mathcal S$
（相应地，$\mathcal S'$）到集合范畴的一个函子，即一个单纯集（相应地，
半单纯集），称为\emph{与范畴 $Z$ 相伴的}单纯集（相应地，半单纯集）；
分别记作 $Z'$ 和 $Z''$。此外，还存在一个显然的自然映射，从 $Z'$ 所给出的
半单纯集到 $Z''$；这个映射当且仅当 $Z$ 是群胚时为同构。由此得到：

\result{命题}{4.1} 从范畴所成的范畴到半单纯集范畴的函子
$Z\longmapsto Z''$ 是全忠实的，并且在\emph{范畴}所成的范畴与下述
半单纯集所成的范畴之间定义一个等价：这些半单纯集即从 $\mathcal S'$ 到
$(\mathrm{Ens})$ 的反变函子 $K$，并且把上述类型的推出
$A\amalg_C B$ 变为集合的纤维积。同样，从群胚范畴到单纯集范畴的函子
$Z\longmapsto Z'$ 是全忠实的，并且在\emph{群胚}范畴与
\src{109}
下述单纯集所成的范畴之间定义一个等价：这些单纯集即从 $\mathcal S$ 到
$(\mathrm{Ens})$ 的反变函子 $K$，并且把推出变为纤维积。

因此，可以把范畴视为特殊的半单纯集，把群胚视为特殊的单纯集；当然，在用
一种结构解释另一种结构时，必须按照惯例“在同构意义下”进行论证。通常的
化归到集合论情形的方法于是推出：

\result{推论}{4.2} 若把前一陈述中的范畴、群胚、单纯集分别换成
$\mathcal C$-范畴、$\mathcal C$-群胚、$\mathcal C$-单纯对象，并假设
$\mathcal C$ 中存在纤维积，则前一陈述仍然成立。

范畴 $\mathcal C$ 中与一个 $\mathcal C$-范畴 $(X,R,\ldots)$ 相伴的
半单纯对象 $K$ 可以明确写出：
可把 $K$ 的分量 $K_n=K(\Delta_n)$ 解释为 $(R,p_1)$ 在 $X$ 上的第
$(n+1)$ 次纤维幂；更好地，也可使用递推公式
\[
  K_0=R,\qquad
  K_n=(K_{n-1},p_n^{(n-1)})\times_X(R,p_1),
\]
\footnote{原印本及法、英可编辑见证均作 $K_0=R$，并在递推式中使用
$p_n^{(n-1)}$；但前文 $K_n=\operatorname{Hom}(\Delta''_n,Z)$ 的定义给出
$K_0=X$、$K_1=R$，且所声明的投影下标范围为
$0\leq i\leq n-1$。这里保留原式；该下标冲突已作为源文疑点记录。}
其中 $p_i^{(n-1)}$（$0\leq i\leq n-1$）是从 $K_{n-1}$ 到 $X$ 的自然
投影（它们本身也由递推定义）。这样，$p_1,p_2,\pi,D,s$ 就分别解释为对应于
$\mathcal S$ 中下列态射的单纯运算：$\Delta_1$ 的面 $0$、
$\Delta_1$ 的面 $1$、$\Delta_2$ 的面 $(0,2)$、退化
$\Delta_1\to\Delta_0$，以及 $\Delta_1$ 的对称。其他所有半单纯运算
（相应地，单纯运算）都可通过复合和纤维积，由前述四个（相应地，五个）
运算形式地导出。

现在，把一个范畴对象 $X$ 上的\emph{预等价关系}定义为如下数据：一个群胚，
其对象之对象（不妨这样说）为 $X$。因此，这种数据尤其包括一个对象 $R$
以及两个态射
\[
  p_1,p_2:R\rightrightarrows X.
\]
不过要注意，与等价对的情形不同，仅凭这些数据并不能确定所考虑的结构。
本报告研究这个概念，是为了得到能够取商的判据，即能够形成态射对
$(p_1,p_2)$ 的余核的判据。因此，问题的陈述并不涉及群胚固有的附加数据。
但在随后结果的证明中，我们仍将
\src{110}
使用这些附加数据，特别是直至维数 $3$ 的单纯运算（其中包括对称运算，并且
涉及 $R$ 在 $X$ 上的四重纤维幂）。

$\mathcal C$ 的对象 $X$ 上的一个等价关系定义一个预等价关系：只需在集合论
情形下验证这一点；对集合 $X$ 上的一个等价关系，取对象集为 $X$、箭头集为
该等价关系之图的群胚即可。

作用在 $\mathcal C$ 的对象 $X$ 上的一个 $\mathcal C$-幺半群 $G$ 定义一个
$\mathcal C$-范畴，其底层对象为 $R=G\times X$ 和 $X$（这里要求
$G\times X$ 存在）；它当且仅当 $G$ 是群时为 $\mathcal C$-群胚。仍只须在
集合论情形下验证这一点。此时，箭头 $(g,a)$ 与 $(g',g\cdot a)$ 的复合定义为
\[
  (g',g\cdot a)(g,a)=(g'g,a),
\]
也就是说，若 $a,b\in X$，则 $\operatorname{Hom}(a,b)$ 按定义是把 $a$
送到 $b$ 的输运子，而态射的复合由 $G$ 中元素的复合给出。

\medskip
\noindent\textbf{注. ---} 对于 (4.1) 这类陈述，只要约定所考虑的全部对象
都位于某个固定的“宇宙”中（该宇宙本身是一个集合），就可以避开它所引起的
逻辑困难。

\fsection{5}{按有限平坦等价关系取商}

\result{定理}{5.1} 设 $X=\operatorname{Spec}(B)$ 是仿射概形，$R$ 是 $X$
中的一个预等价关系，其分量 $R_1$ 是仿射的，记
$R_1=\operatorname{Spec}(C)$。假设第一投影 $p_1:R_1\to X$ 是有限局部自由
态射；换言之，相应的环同态 $p_1^*:B\to C$ 使 $C$ 成为有限生成投射
$B$-模。令 $A$ 为 $B$ 的如下子环：它是同态对
\[
  p_1^*,p_2^*:B\rightrightarrows C
\]
的核（即所有满足 $p_1^*(b)=p_2^*(b)$ 的元素 $b$ 所成的集合）。令
$Y=\operatorname{Spec}(A)$，并令 $f:X\to Y$ 为由 $A$ 到 $B$ 的嵌入所定义
的态射。在这些条件下：
\begin{enumerate}
\item[(i)] $B$ 在 $A$ 上整，即 $f$ 是整态射。

\item[(ii)] 态射 $f$ 是满射，其纤维是下述等价类
\src{111}
$p_2\allowbreak(p_1^{-1}(x))$，即 $X$ 模 $R$ 的集合论等价类；而且 $Y$
的拓扑是 $X$ 的拓扑的商。

\item[(iii)] $Y$ 是预概形范畴中 $X$ 按 $R$ 所取的商。

\item[(iv)] 若 $R$ 来自一个\emph{等价关系}，则态射 $f:X\to Y$ 是有限
局部自由的（即 $B$ 是有限生成投射 $A$-模），且该等价关系是有效的，即
$R_1\to X\times_Y X$ 是同构。
\end{enumerate}

这个定理推广了如下熟知定理：有限群 $G$ 以自同构作用于环 $B$，而 $A$ 是
不变量环；其证明也与熟知的证明类似。可将 (iii) 具体化如下：

\result{推论}{5.2} 典范态射
$R_1\to X\times_Y X$ 是\emph{满射}。

仍设 $R$ 是 $X$ 中一个“有限且局部自由”的预等价关系，但现在允许 $X$ 为
任意预概形。假设能够找到一个预概形 $Y$ 和一个态射 $f:X\to Y$，使得
$fp_1=fp_2$；并且还假设 $Y$ 上环层同态序列
\[
  \mathcal O_Y\longrightarrow f_*(\mathcal O_X)
  \rightrightarrows g_*(\mathcal O_R)
\]
正合（其中 $g=fp_i$）。由上述定理可推出与定理中 (i) 至 (iv) 类似的结论；
特别地，由 (iii)，$Y$ 是 $X$ 按 $R$ 所取的商，因而在唯一同构意义下确定。
在这些条件下，称 $X$ 中的预等价关系 $R$ 是\emph{可容许的}。采用这一定义：

\result{定理}{5.3} 设 $X$ 是预概形，$R$ 是 $X$ 中的预等价关系，并且
$p_1:R_1\to X$ 是有限局部自由态射。则 $R$ 可容许，当且仅当 $X$ 模 $R$
的每个集合论等价类 $p_2(p_1^{-1}(x))$ 都包含在某个仿射开集内（若 $X$
的每个有限子集都包含在某个仿射开集内，则这一条件总是成立；例如，$X$ 在
某个仿射概形上拟射影时即如此）。

事实上，不难证明 $X$ 模 $R$ 的每个等价类都包含在某个\emph{在 $R$ 下稳定}
的仿射开集内；再把利用定理 5.1 得到的各个部分粘合起来，便构造出商 $Y$。
\src{112}

\result{推论}{5.4} 假设上述条件成立，并进一步假设 $R$ 来自一个等价关系。
那么，该等价关系是有效的，即 $R_1\to X\times_Y X$ 是同构，而且
$f:X\to Y$ 是有限局部自由态射。

由“下降”立即推出：

\result{推论}{5.5} 在 (5.4) 的条件下，$X$ 在 $Y$ 上处处秩为 $n$，当且
仅当 $(R_1,p_1)$ 在 $X$ 上处处秩为 $n$。若 $X,R_1$ 是 $Z$-预概形，
$p_1,p_2$ 是 $Z$-态射，因而 $Y$ 是 $Z$-预概形，则 $X$ 在 $Z$ 上平坦
当且仅当 $Y$ 在 $Z$ 上平坦。

综上：

\result{附论}{} 给出预概形之间的一个有限、局部自由且满的态射
$f:X\to Y$，等价于给出一个配有等价关系 $R$ 的预概形 $X$，使得
$p_1:R\to X$ 有限且局部自由，并且每个等价类
$p_2(p_1^{-1}(x))$ 都包含在某个仿射开集内。

\medskip
\noindent\textbf{注 5.6.}

\smallskip
\noindent $1^\circ$ 以上没有使用任何诺特性假设。

\smallskip
\noindent $2^\circ$ 这种取商概念把 CARTIER 的“不可分下降”包括为特殊情形；
后者对应于确定如下有限局部自由态射 $f:X\to Y$：$f_*(\mathcal O_X)$
具有相对于 $\mathcal O_Y$ 的 $p$-基（其中 $X$ 是给定的预概形，且其层
$\mathcal O_X$ 被素数 $p>0$ 零化）。注意，这一结果很容易在不对局部环施加
正则性假设、也不假设 $X$ 是域上的代数概形的情况下表述。取 $X$ 为特征
$p$ 的域的谱，便得到 JACOBSON--BOURBAKI 理论。

\smallskip
\noindent $3^\circ$ GABRIEL 先前已在域 $k$ 上有限交换群的取商理论中得到了
定理 (5.3) 的一个特殊情形。[参见 (7.4)]。

\fsection{6}{按固有平坦等价关系取商}

\result{定理}{6.1} 设 $S$ 是局部诺特预概形，$X$ 是拟射影 $S$-概形，
$R$ 是 $X$ 中的一个预等价关系，并且满足：
\src{113}
\begin{enumerate}
\item[(a)] $p_1:R_1\to X$ 是固有且平坦的；
\item[(b)] $R_1\to X\times_S X$ 是有限态射；或者，由 (a)，等价地说，
它的纤维是有限的（若 $R$ 来自一个等价关系，则此条件自动满足）。
\end{enumerate}
在这些条件下：
\begin{enumerate}
\item[(i)] $X/R=Y$ 存在，并且（若 $S$ 是诺特的且 $R$ 来自一个等价关系）
在 $S$ 上拟射影。\footnote{印刷文本没有第二个假设。与报告 212 相联系的勘误
指出，此时似乎只有在 $R$ 来自一个等价关系时，拟射影性才得到了证明。}

\item[(ii)] 典范态射 $f:X\to Y$ 是满射、固有且开的；它的纤维是 $X$
模 $R$ 的等价类 $p_2(p_1^{-1}(x))$。因此，$Y$ 可与 $X$ 关于 $R$ 所定义
之集合论等价关系的商拓扑空间等同。最后，
$R_1\to X\times_Y X$ 是满射。

\item[(iii)] 若 $R$ 来自一个等价关系，则该等价关系是有效的，即
$R_1\to X\times_Y X$ 是同构；而且 $f:X\to Y$ 是平坦的
（从而是忠实平坦的）。
\end{enumerate}

在证明中，借助 $X$ 关于 $R$ 的适当准截面可将问题化归到 (5.1)；
其证明类似于 Chevalley 讨论班中商代数群的构造。

概括如下：

\result{附论}{} 设 $X$ 在局部诺特的 $S$ 上拟射影。给定一个从 $X$ 到
某个 $S$-预概形 $Y$ 的固有、忠实平坦且满的态射 $f:X\to Y$，等价于
给定 $X$ 上的一个等价关系 $R$，使得 $p_1:R\to X$ 固有且平坦。

同一方法给出下列结果：

\result{定理}{6.2} 设 $S$ 是诺特预概形，$X$ 是 $S$ 上有限型的预概形，
$R$ 是 $S$-预概形 $X$ 中的一个预等价关系。假设：
\begin{enumerate}
\item[(a)] $p_1:R_1\to X$ 平坦且为有限型；
\item[(b)] 态射 $R_1\to X\times_S X$ 是拟有限的（即纤维有限）。
\end{enumerate}
则 $X$ 中存在一个关于 $R$ \emph{饱和}的稠密开子集 $U$，使得：
\begin{enumerate}
\item[(i)] 若 $R_U$ 是 $R$ 在 $U$ 中诱导的预等价关系，则 $U/R_U$
存在，并且在 $S$ 上为有限型。

\item[(ii)] 典范态射 $U\to U/R_U$ 是满射且为开态射；它的纤维是
$R_U$ 的集合论等价类（因而 $U/R_U$ 是 $U$ 关于 $R_U$ 所定义之
集合论等价关系的商拓扑空间）。最后，态射
\[
  (R_U)_1\longrightarrow U\times_{U/R_U}U
\]
是满射。
\src{114}

\item[(iii)] 若 $R$ 来自一个等价关系，则可以使
$U\to U/R_U$ 忠实平坦，并使 $R_U$ 有效。
\end{enumerate}

这是一个本质上具有“双有理”性质的结果。

\medskip
\noindent\textbf{评注 6.3。}

\smallskip
\noindent $1^\circ$ 我不知道在 (6.1) 和 (6.2) 中假设 (b) 是否多余。
在取群的商时，它实际上迫使我们限于所有稳定子都是有限群的情形。

\smallskip
\noindent $2^\circ$ 可以问，在没有平坦性假设时，是否仍有与 (6.1)
和 (6.2) 类似的结果。我不知道这方面的任何反例。另一方面，即使保留
平坦性假设，并且只考虑满足 $p_1:R\to X$ 平坦且拟有限（但不有限）
的等价关系，当 $X$ 仿射时，$R$ 仍可能不是有效的：取覆盖 Nagata 簇的
仿射开集上所诱导的等价关系即可（在该簇上，一个二阶群以“非容许”的方式
作用）。

\fsection{7}{应用}

正如引言中所说，(6.1) 最重要的应用是构造 Picard 概形，以及求解其他
各种“模问题”；我们以后还会回到这些问题。

由此可对 SHIMURA 的下列结果给出一个简单证明：

\result{命题}{7.1} 设 $A$ 是定义在离散赋值环 $V$ 上的阿贝尔概形，
$K$ 是 $V$ 的分式域。那么，$K$ 上每个与 $A\otimes_V K$ 的某个商同源的
阿贝尔概形 $B'$ 都“在 $V$ 上具有良好约化”，即同构于某个
$B\otimes_V K$，\footnote{印刷文本作 $B\otimes_V A$。由于 $B$ 是
$V$-概形，而命题所讨论的是它的泛纤维，基因子必为 $K$。}其中 $B$ 是
$V$ 上的阿贝尔概形（回顾一下，它在本质上是唯一的）。

可以设 $B'$ 是 $A_K$ 关于某个子群概形 $C'$ 的商。（注意：$C'$ 一般
并不是“既约”的，即它的局部环中会有幂零元。）考虑 $A$ 的闭子概形 $C$，
它是 $C'$ 的“闭包”，即满足 $C_K$ 包含 $C'$ 的最小闭子概形。将有
$C_K=C'$；由 $V$ 是离散赋值环，很容易推出 $C$ 是 $V$ 上 $A$ 的
子群概形。由于 $A$ 在 $\operatorname{Spec}(V)=S$ 上固有，$C$ 也如此；
另一方面，$A$ 在 $S$ 上甚至是射影的。因此可以应用 (6.1) 构造
$A/C=B$，
\src{115}
这正是所求的 $B$。

最后，利用本质上已知的论证，可以从 (6.2) 推出下列结果：

\result{定理}{7.2} 设 $S$ 是一个阿廷环的谱，$F$ 和 $G$ 是 $S$ 上
有限型的群概形，$u:F\to G$ 是 $S$ 上群概形的同态。假设：
\begin{enumerate}
\item[(i)] $F$ 在 $S$ 上平坦；
\item[(ii)] $u$ 的核是有限的。
\end{enumerate}
在这些条件下，商概形 $G/F$ 存在；典范态射 $G\to G/F$ 是满射且为
开态射，其纤维是由 $F$ 在 $G$ 上的右作用所定义的集合论等价类。最后，
若 $u$ 是单态射，则态射 $G\to G/F$ 平坦，并且态射
\[
  G\times F\longrightarrow G\times_{G/F}G
\]
是同构。换言之，$G$ 是 $G/F$ 上的主齐性空间，其结构群（从右作用）为
$F$；更确切地说，是把 $F\times_S(G/F)$ 看作 $G/F$ 上的群概形
（参见 [1]，B，\S~6）。

\result{推论}{7.3} 在这些条件下，$G$ 在 $S$ 上平坦，当且仅当 $G/F$
在 $S$ 上平坦。若此条件成立，则关于 $F$ 取商与 $S$ 的任意基扩张交换；
若 $F$ 是 $G$ 的正规子群，则可在 $G/F$ 上赋予 $G$ 关于 $F$ 的
\emph{商群}结构。

当 $S$ 是一个域的谱时，情形尤其简单，因为每个 $S$-预概形都自动在
$S$ 上平坦。由此得到：

\result{推论}{7.4} 设 $F,G$ 是域 $k$ 上有限型的群概形，
$u:F\to G$ 是 $k$-群同态。则 $u$ 可分解为 $F\to F'\to G$，其中
$F\to F'$ 是关于 $F$ 的闭子群 $\operatorname{Ker}u$ 取商的同态，
而 $F'\to G$ 是一个闭浸入的群同态。商 $G/F=G/F'$ 存在。在 $k$ 上的
代数群中，通常的形式体系（如诺特型定理）成立。

这一结果使我们能够统一处理经典意义下代数群的取商（即在 $k$ 上不可约且
在 $k$ 上光滑的代数群），以及关于 CARTIER 所考察的“无穷小”子群
\src{116}
取商。把 CARTIER 在 DIEUDONNÉ 关于形式群的工作之后引入的“超代数”
看作 $k$ 上形式概形范畴中的群，是有利的；在适用时（即它们对应于
$k$ 上有限秩的超代数时），还可把它们看作 $k$ 上的有限代数群。

\fsection{8}{一个猜想}

这一猜想源于如下需要：对作用在“Hilbert 概形”的某些子概形上的射影群
取商（在概形理论中，Hilbert 概形取代了 Chow 簇）。

设 $S$ 是预概形，$n$ 是整数。对 $S$ 上的每个预概形 $S'$，令其对应于
群 $\operatorname{Gl}(n,\Gamma(S',\mathcal O_{S'}))$，即元素取自
$\mathcal O_{S'}$ 的截面环的 $n$ 行 $n$ 列可逆矩阵所成的群。由此得到
关于 $S'$ 的反变函子；容易证明它是可表的，因而对应于 $S$ 上的一个群概形，
而且该群概形在 $S$ 上仿射，记作 $\operatorname{Gl}(n)_S$。它的构造与
基变换相容，因此实际上这一切都来自 $\mathbf Z$ 上的一个群概形，记作
$\operatorname{Gl}(n)$。群 $\operatorname{Gl}(1)$ 称为乘法群，常记作
$\mathbf G_m$；它对应于函子
$S\longmapsto\Gamma(S,\mathcal O_S)^*$，即 $S$ 上“单位”所成的群。
有一个显然的同态 $\operatorname{Gl}(1)\to\operatorname{Gl}(n)$；
容易构造其商群，记作 $\operatorname{GP}(n-1)$，称为
\emph{$\mathbf Z$ 上次数为 $n-1$ 的射影群}。它表示如下函子：把 $S$
映到层
$\widetilde{\operatorname{Gl}(n)}_S/
\widetilde{\operatorname{Gl}(1)}_S$ 的截面群，其中
$\widetilde{\operatorname{Gl}(n)}_S$ 表示
$\operatorname{Gl}(n)_S$ 在 $S$ 上的截面芽所成的层。（须注意：
$\operatorname{GP}(n-1)_S$ 在 $S$ 上的截面一般并非来自
$\operatorname{Gl}(n)_S$ 在 $S$ 上的截面！）还应指出，至少在 $S$
诺特时，可以证明 $\operatorname{GP}(n-1)$ 也表示函子
$S\longmapsto\operatorname{Aut}_S(\mathbf P_S^{n-1})$，其中
$\mathbf P_S^{n-1}$ 是 $S$ 上相对维数为 $n-1$ 的标准射影概形。
它正是以这种方式进入模理论的。

设 $S$ 是诺特概形；如果愿意，可以假设它是仿射的。设 $X$ 是拟射影
$S$-预概形，并配备一个相对于 $S$ 甚丰沛的可逆层 $\mathcal L$。
假设群 $G=\operatorname{GP}(n)_S$ 作用在 $X$ 上，同时也作用在
$\mathcal L$ 上（并与它在 $X$ 上的作用相容），而且它在 $X$ 上
\emph{自由地}作用。\footnote{印刷文本作“在 $S$ 上”。这里所考察的
作用和取商条件都是关于 $X$ 的。}
\src{117}

\result{猜想}{8.1}\footnote{与报告 212 相联系的勘误指出，这一猜想
“终究是错误的”。它还指出，评注 8.1 中的肯定性结果后来已经得到证明，
并提及 Mumford 的工作及其改进。} 在上述条件下：

\smallskip
\noindent $1^\circ$ \emph{$G$ 所定义的等价关系是有效的，商
$Y=X/G$ 在 $S$ 上为有限型，并且典范态射 $f:X\to Y$ 平坦且为满射}
（因而 $X$ 成为 $Y$ 上的主齐性丛，其群为
$G\times_S Y=\operatorname{GP}(n)_Y$）。

\smallskip
\noindent $2^\circ$ 设 $\mathcal L'$ 是由 $\mathcal L$ 沿 $f$ 作
“忠实平坦下降”而得到的 $Y$ 上可逆层（参见 [1]，B，\S~定理 1）。
则 $\mathcal L'$ \emph{相对于 $S$ 在 $Y$ 上预丰沛}，即存在一个整数
$m$ 和一个从 $Y$ 到某个适当的标准射影概形 $\mathbf P_S^N$ 的拟有限态射，
使 $(\mathcal L')^{\otimes m}$ 同构于
$\mathcal O_{\mathbf P_S^N}(1)$ 的逆像。

应当指出，即使 $X$ 在 $S$ 上是分离的，$Y$ 仍可能不在 $S$ 上分离
（这种情形会出现在完全不病态的“模问题”中）。若 $1^\circ$ 成立，则
$Y$ 分离，当且仅当 $G$ 所定义的等价关系具有闭图，即
$G\times X\to X\times X$ 的像是闭的（此时该态射是闭浸入）。
若 $Y$ 分离，则 $\mathcal L'$ 相对于 $S$ 在 $Y$ 上预丰沛，当且仅当
它是丰沛的，即某个适当的张量幂给出一个射影浸入。在引言所述的模问题中，
可以证明最终得到的等价关系确实具有闭图。

\medskip
\noindent\textbf{评注 8.1。---} 为了使讨论具体，并且因为目前它在实践中
是最重要的情形，我们假设了 $G=\operatorname{GP}(n)_S$。对 $G$ 更合理的
假设似乎应当是：$G$ 是某个 Tohokû 群在 $S$ 上的一种“形式”
（CHEVALLEY 也已在整数上构造了这些群）。就上述猜想而言，我所知道的
唯一肯定性结果如下：设 $X$ 是\emph{特征为 $0$ 的域 $k$ 上的仿射概形，
群 $\operatorname{Gl}(n)_k$ 或 $\operatorname{GP}(n-1)_k$ 在其上
自由地作用。则 $G$ 所定义的等价关系是有效的，商 $X/G$ 也是仿射的，
并且态射 $X\to X/G$ 平坦且为满射}。证明使用下述事实（目前仅在特征
$0$ 时得到证明）：让 $G$ 作用在 $G$ 的仿射坐标环 $A$ 上，并把 $A$
看作 $k$ 上的向量空间，则 $G$ 的平凡表示只出现一次（在某个 $k$ 上
有限维且在 $G$ 下稳定的向量子空间的合成列中）。看来，系统地运用
\src{118}
$G$ 的线性表示理论，最终可能给出这一猜想的证明，至少当基底是一个域时
如此。当基底不再是域时，报告人对此一无所知。

\section*{参考文献}
\addcontentsline{toc}{section}{参考文献}

\begin{description}
\item[{[1]}] GROTHENDIECK (Alexander). --- 下降技术与代数几何中的
存在性定理，I：一般论，Séminaire Bourbaki，1959/60，
n\textsuperscript{o} 190，29 页。

\item[{[2]}] GROTHENDIECK (Alexander). --- 下降技术与代数几何中的
存在性定理，II：形式模理论中的存在性定理，Séminaire Bourbaki，
1959/60，n\textsuperscript{o} 195，22 页。

\item[{[3]}] GROTHENDIECK (A.) 与 DIEUDONNÉ (J.). --- 代数几何基础，
I：概形的语言。--- 巴黎，Presses universitaires de France，1960
（Institut des hautes Études scientifiques, Publications
mathématiques, 4）。
\end{description}
